\begin{table*}[t]
    \centering
    \small
    \renewcommand{\arraystretch}{1.15}
    \begin{tabular}{p{0.17\textwidth} p{0.77\textwidth}}
        \toprule
        \textbf{Symbol} & \textbf{Meaning} \\
        \midrule
        \multicolumn{2}{@{}l}{\textit{Time series \& sampling}} \\[-1pt]\cmidrule{1-2}
        $T$           & Total observation duration. \\
        $\Delta_t$    & Sampling interval; $\Delta_t = 1/f_s$. \\
        $f_s$         & Sampling frequency. \\
        $f_{Ny}$      & Nyquist frequency, $f_{Ny} = f_s/2$. \\
        $n$           & Number of time samples per channel, $n = T/\Delta_t$. \\
        $p$           & Number of channels. \\
        $\Delta_f$    & Frequency resolution, $\Delta_f = 1/T$. \\
        $f_k$         & $k$th discrete Fourier frequency, $f_k = k\Delta_f = k/T$. \\
        $N$           & Number of positive Fourier frequencies: $N=n/2$ ($n$ even) or $(n-1)/2$ ($n$ odd). \\
        \multicolumn{2}{@{}l}{\textit{Blocking \& coarse-graining}} \\[-1pt]\cmidrule{1-2}
        $N_b$         & Number of non-overlapping time blocks. \\
        $L_b$         & Samples per block, $L_b = n/N_b$. \\
        $T_b$         & Block duration, $T_b = L_b\Delta_t = T/N_b$. \\
        $N_l$         & Positive-frequency count on the fine (per-block) grid, $N_l = L_b/2$. \\
        $N_c$         & Number of coarse-grain frequency bins. \\
        $N_h$         & Frequencies aggregated into each coarse bin, $|J_h|$. \\
        $J_h$         & $h$th coarse-grain bin: a consecutive subset of fine-grid frequencies. \\
        $\bar{f}_h$   & Midpoint frequency of bin $J_h$. \\
        \multicolumn{2}{@{}l}{\textit{Spectral density}} \\[-1pt]\cmidrule{1-2}
        $\Z$          & Observed $p$-channel time series, $\Z=(\Z_1,\ldots,\Z_n)^\intercal\in\mathbb{R}^{n\times p}$. \\
        $\d(f_k)$     & DFT vector at $f_k$; $\d(f_k)=(d_1(f_k),\ldots,d_p(f_k))^\intercal\in\mathbb{C}^p$. \\
        $\Gamma(h)$   & Autocovariance matrix at lag $h$, $\Gamma(h)=(\gamma_{lm}(h))=\mathbb{E}(\Z_t\Z_{t+h}^\intercal)$. \\
        $\gamma_{lm}(h)$ & $(l,m)$ entry of $\Gamma(h)$: cross-covariance between channels $l$ and $m$ at lag $h$. \\
        $\S(f_k)$     & $p\times p$ Hermitian positive-definite spectral density matrix at $f_k$; Fourier transform of $\Gamma(h)$. \\
        \multicolumn{2}{@{}l}{\textit{Periodogram \& Wishart statistics}} \\[-1pt]\cmidrule{1-2}
        $\I^{(i)}(f_k)$ & Periodogram matrix of block $i$: $\I^{(i)}(f_k)=\d^{(i)}(f_k)\d^{(i)}(f_k)^*\sim\mathcal{CW}_p(T_b\S(f_k),1)$. \\
        $\bar{\I}(f_k)$ & Block-averaged periodogram, $\bar{\I}(f_k)=\frac{1}{N_b}\sum_{i=1}^{N_b}\I^{(i)}(f_k)$. \\
        $\Y(f_k)$     & Summed periodogram / Wishart statistic, $\Y(f_k)=N_b\bar{\I}(f_k)\sim\mathcal{CW}_p(T_b\S(f_k),N_b)$. \\
        $\Y_h^{*}$    & Coarse-grained statistic for bin $J_h$: $\Y_h^*=\sum_{f_k\in J_h}\Y(f_k)\,\dot{\sim}\,\mathcal{CW}_p(T_b\S(\bar{f}_h),N_bN_h)$. \\
        $\lambda_\nu^{(k)}$ & $\nu$th eigenvalue of $\Y(f_k)$; used to define scaled eigenvectors $\u_\nu^{(k)}$. \\
        $\v_\nu^{(k)}$      & $\nu$th unit eigenvector of $\Y(f_k)$, so $\Y(f_k)=\sum_\nu\lambda_\nu^{(k)}\v_\nu^{(k)}\v_\nu^{(k)*}$. \\
        $\u_\nu^{(k)}$      & Scaled eigenvector, $\u_\nu^{(k)}=\sqrt{\lambda_\nu^{(k)}}\,\v_\nu^{(k)}$, so $\Y(f_k)=\sum_\nu\u_\nu^{(k)}\u_\nu^{(k)*}$. \\
        $\U^{(k)}$    & $p\times p$ matrix with columns $\u_\nu^{(k)}$; rows indexed as $\u_{(j)}^{(k)}$. \\
        \multicolumn{2}{@{}l}{\textit{Cholesky parametrisation}} \\[-1pt]\cmidrule{1-2}
        $\mathbf{T}_k$ & Unit lower-triangular matrix in the Cholesky factor: $\S(f_k)^{-1}=\mathbf{T}_k^*\mathbf{D}_k^{-1}\mathbf{T}_k$. \\
        $\mathbf{D}_k$ & Diagonal matrix with positive entries $\delta_{1k}^2,\ldots,\delta_{pk}^2$. \\
        $\theta_{jl}^{(k)}$ & Complex-valued $(j,l)$ off-diagonal element of $\mathbf{T}_k$ (for $j>l$), at frequency $f_k$. \\
        $\delta_{jk}^2$     & $j$th diagonal element of $\mathbf{D}_k$ at frequency $f_k$; controls the $j$th partial variance. \\
        $\btheta_j$   & Vector of all $\theta_{jl}^{(k)}$ for component $j$, across frequencies and $l<j$. \\
        $\bdelta_j$   & Vector of all $\delta_{jk}$ for component $j$, across frequencies $k=1,\ldots,N_l$. \\
        $r_{j\nu}^{(k)}$ & Cholesky residual for component $j$: $r_{j\nu}^{(k)} := u_{j\nu}^{(k)} - \sum_{l=1}^{j-1}\theta_{jl}^{(k)}u_{l\nu}^{(k)}$. \\
        \multicolumn{2}{@{}l}{\textit{P-spline model}} \\[-1pt]\cmidrule{1-2}
        $K$           & Number of B-spline basis functions (basis size); note $K$ is a model hyperparameter, distinct from the frequency index $k$. \\
        $B_m(f_k)$   & $m$th B-spline basis function evaluated at $f_k$, $m=1,\ldots,K$. \\
        $M_j$, $M_{jl}$ & Number of basis functions used for the $j$th diagonal ($M_j$) and $(j,l)$ off-diagonal ($M_{jl}$) components. \\
        $\bold{w}_j^{(\delta)}$ & Spline coefficients for $\log\delta_{jk}^2$; $\log\delta_{jk}^2 = \sum_m B_m(f_k)\,w_{j,m}^{(\delta)}$. \\
        $\bold{w}_{jl}^{(\Re)}$, $\bold{w}_{jl}^{(\Im)}$ & Spline coefficients for the real and imaginary parts of $\theta_{jl}(f_k)$. \\
        $\bold{P}_j$  & Penalty matrix for component $j$: $[\bold{P}_j]_{lm}=\int_0^1 B_l''(t)\,B_m''(t)\,\mathrm{d}t$. \\
        $\phi_j$      & Smoothing precision hyperparameter for component $j$; controls strength of the spline penalty. \\
        \bottomrule
    \end{tabular}
    \caption{Notation used throughout the paper. Symbols are grouped by the stage of the model they belong to.}
    \label{tab:definitions}
\end{table*}
